\documentclass[11pt,a4paper]{amsart}
\usepackage[T1]{fontenc}
\usepackage{lmodern}
\usepackage{amssymb}
\usepackage{microtype}
\usepackage[hyphens]{url}
\usepackage[colorlinks=true,linkcolor=black,citecolor=black,urlcolor=black]{hyperref}
\usepackage[margin=2.5cm]{geometry}

\hypersetup{
  pdftitle={Which real quadratic planes need four colours},
  pdfauthor={Sergi Gal\'an},
  pdfsubject={The Hadwiger-Nelson problem over real quadratic fields},
  pdfkeywords={Hadwiger-Nelson problem, chromatic number of the plane, unit-distance graph, real quadratic field,
  Cayley graph, character}}

\newtheorem{theorem}{Theorem}
\newtheorem{corollary}[theorem]{Corollary}
\newtheorem{proposition}[theorem]{Proposition}
\newtheorem{lemma}[theorem]{Lemma}
\theoremstyle{remark}
\newtheorem*{remark}{Remark}

\newcommand{\Q}{\mathbb{Q}}
\newcommand{\R}{\mathbb{R}}
\newcommand{\Z}{\mathbb{Z}}
\newcommand{\C}{\mathbb{C}}
\newcommand{\F}{\mathbb{F}}
\newcommand{\Cay}{\operatorname{Cay}}
\newcommand{\Sq}{\mathrm{Sq}}

\title{Which real quadratic planes need four colours}
\author{Sergi Gal\'an}
\thanks{These results were obtained with the assistance of an AI system (Claude, by Anthropic), which found the
construction, wrote the proofs and the programs that check them, and helped to draft this note. Separate AI agents
refereed the proofs, with their own programs. No mathematician has checked the note yet. The author is not a
professional mathematician and submits it for independent verification.}
\date{Draft, 3 October 2026}
\subjclass[2020]{Primary 05C15, 52C10; Secondary 11R11, 05C25}
\keywords{Hadwiger--Nelson problem, chromatic number of the plane, unit-distance graph, real quadratic field,
Cayley graph, character}

\begin{document}

\begin{abstract}
For a subfield $F\subset\R$ let $\chi(F^2)$ be the chromatic number of the graph on $F^2$ that joins two points at
Euclidean distance $1$. For a real quadratic field $\Q(\sqrt d)$, $d$ squarefree, it is known that
$\chi(\Q(\sqrt d)^2)=2$ unless $d\equiv 3\pmod 4$, that $\chi(\Q(\sqrt d)^2)\ge3$ if $d\equiv3\pmod 4$, and that
$\chi(\Q(\sqrt d)^2)\le 3$ unless $d\equiv 11\pmod{12}$. We prove that $\chi(\Q(\sqrt d)^2)\ge 4$ for every
$d\equiv 11\pmod{12}$. So the plane over $\Q(\sqrt d)$ needs four colours exactly when $d\equiv11\pmod{12}$, and,
with the upper bounds of Fischer and Moorhouse, $\chi(\Q(\sqrt d)^2)=4$ for all these $d$ except possibly
$d\equiv 47,143,167\pmod{168}$. The proof exhibits, for each $d$, a finite set of unit vectors, made of $(1,0)$, a
vector and its mirror image, one more vector when $d\equiv11\pmod{24}$, and their rotations by the rational rotations
whose denominators divide $5^k$ for a suitable $k$, and shows that no
character of the group they generate maps all of them into $[1/3,2/3]$; a characterisation of $3$-colourable
Cayley graphs of abelian groups by such characters then excludes $3$-colourings. The main step is an exact
description of the characters of $\frac1N\Z[i]$, $N=5^k$, that keep every rational rotation with denominator
dividing $N$ in $[1/3,2/3]$.
\end{abstract}

\maketitle

\section{Introduction}

Let $F\subset\R$ be a field. The \emph{plane over $F$} is the graph on $F^2$ in which $x$ and $y$ are adjacent when
$|x-y|=1$, and $\chi(F^2)$ is its chromatic number. For $F=\R$ this is the Hadwiger--Nelson problem, and
$6\le\chi(\R^2)\le7$: the bound $\chi(\R^2)\ge5$ is due to de Grey \cite{deGrey} and to Exoo and Ismailescu
\cite{EI}, the bound $\chi(\R^2)\ge6$ to a recent non-constructive proof \cite{OpenAI}, and the upper bound to
Isbell (see \cite{Soifer}). For $F=\Q$, $\chi(\Q^2)=2$ \cite{Woodall}. For real quadratic fields
$K=\Q(\sqrt d)$, $d\ge2$ squarefree, the following was known:
\begin{itemize}
\item $\chi(K^2)=2$ if $d\not\equiv3\pmod4$ \cite{Johnson1987,Fischer1990,Moorhouse};
\item $\chi(K^2)\ge3$ if $d\equiv3\pmod4$ \cite[Thm.~8]{Fischer1990} (and \cite[Thm.~8.6]{Moorhouse} when $K$
contains $\sqrt p$ for a prime $p\equiv3\pmod4$); for
$u=((1-d)/(1+d),\,2\sqrt d/(1+d))$ and its mirror image $\bar u$, the relation
$\frac{d+1}4(u+\bar u)+\frac{d-1}2(1,0)=0$ is a closed walk of odd length $d$;
\item $\chi(K^2)\le3$ if $d\not\equiv2\pmod3$ \cite[Thm.~9]{Fischer1990}, \cite[Cor.~8.3]{Moorhouse};
\item $\chi(K^2)\le4$ if $d\equiv3\pmod8$ \cite[Thm.~10]{Fischer1990} or $d\equiv0,1,2,4\pmod7$
\cite[Cor.~8.3]{Moorhouse}; Moorhouse's Theorem~8.1 combines his bounds at $3$ and $7$ with his Lemma~8.4:
$\chi(K^2)\le4$ unless $d\equiv47,59,83\pmod{84}$.
\end{itemize}
So a real quadratic field can need four colours only when $d\equiv11\pmod{12}$. In the literature we found, the
largest lower bound for a real quadratic field is $3$ \cite{Madore,Moorhouse,Fischer1990}; explicit graphs for $27$
fields that need four colours are given in the companion draft \cite{QP}. Moorhouse
asked what can be said about $\chi(\Q(\sqrt d)^2)$, and about $\chi(\Q(\sqrt{47})^2)$ in particular
\cite{MoorhouseTalk}. Fischer \cite[Thm.~10(ii)]{Fischer1990} showed, with the vector $u$ above (up to sign and
mirror image), that for $d\equiv-1\pmod{24}$ the vectors $(\frac12,0)$ and $(\frac13,0)$ are sums of unit vectors,
so that the plane has no \emph{additive} $k$-colouring for $k\le6$: none induced by a homomorphism from the group
$C_0$ generated by the unit vectors to a group of order $k$ (his Theorem~1 extends colourings of $C_0$ to the plane by
cosets).

\begin{theorem}\label{thm:main}
$\chi(\Q(\sqrt d)^2)\ge4$ for every positive integer $d\equiv11\pmod{12}$.
\end{theorem}

\begin{corollary}\label{cor:classification}
For squarefree $d\ge2$, $\chi(\Q(\sqrt d)^2)\ge4$ if and only if $d\equiv11\pmod{12}$. For these $d$,
$\chi(\Q(\sqrt d)^2)=4$ unless $d\equiv47,143,167\pmod{168}$; in particular $4\le\chi(\Q(\sqrt{47})^2)\le5$.
\end{corollary}

\begin{proof}
The first statement is Theorem~\ref{thm:main} and the bounds listed above. Among $d\equiv11\pmod{12}$, those not
covered by Fischer's or Moorhouse's bound $4$ satisfy $d\equiv7\pmod8$ and $d\equiv3,5,6\pmod7$, that is,
$d\equiv47,143,167\pmod{168}$. For $d=47$, the bound $5$ comes from reduction at a prime above $11$, which splits
in $\Q(\sqrt{47})$, and $\chi(\F_{11}^2)=5$ \cite[Lemma~8.2 and Table~6.1]{Moorhouse}.
\end{proof}

\begin{corollary}\label{cor:padic}
For every prime $p\ge5$, $\chi(\Q_p^2)\ge4$, where $\Q_p^2$ is the graph joining $x,y\in\Q_p^2$ when
$(x_1-y_1)^2+(x_2-y_2)^2=1$. So $\chi(\Q_p^2)\ge4$ exactly for $p\ge5$.
\end{corollary}

\begin{proof}
As $\gcd(p,12)=1$, some positive integer $d$ satisfies $d\equiv11\pmod{12}$ and $d\equiv1\pmod p$ (it need not be
squarefree; Theorem~\ref{thm:main} covers it). By Hensel's lemma $d$ has a square root in $\Z_p^\times$, so
$\sqrt d\mapsto s$ defines a field embedding $\sigma\colon\Q(\sqrt d)\to\Q_p$. Being an injective ring homomorphism
fixing $\Q$, $\sigma$ preserves the relation $(x_1-y_1)^2+(x_2-y_2)^2=1$ in both directions, so applied coordinatewise
it maps the plane over $\Q(\sqrt d)$ isomorphically onto an induced subgraph of $\Q_p^2$, and
$\chi(\Q_p^2)\ge\chi(\Q(\sqrt d)^2)\ge4$. For $p=2,3$ no $d\equiv11\pmod{12}$ is a square in $\Q_p$, and indeed
$\chi(\Q_2^2)=2$ and $\chi(\Q_3^2)=3$, following from \cite{Madore}.
\end{proof}

For $p\equiv1\pmod4$ this is weaker than $\chi(\Q_p^2)=\infty$, a consequence of a theorem of Davies \cite{Davies}; the
new case is $p\equiv3\pmod4$, where explicit graphs over $27$ fields had given the bound for every such $p$ with
$7\le p<2\,129\,503\,819$ but not for $p=2\,129\,503\,819$, and
where no finite set of fields can give it for every $p$, by Chebotarev's theorem \cite{PA}. A finite $4$-chromatic
unit-distance graph in $\Q_p^2$ exists, by compactness, but is not exhibited.

The proof is explicit. Let $G_N$ be the set of the $4(2k+1)$ rational unit vectors whose denominators divide
$N=5^k$ (Lemma~\ref{lem:GN}), and for an integer $n$ let
\[
u_n=\frac{(n+i\sqrt d)^2}{n^2+d}=\Bigl(\frac{n^2-d}{n^2+d},\ \frac{2n\sqrt d}{n^2+d}\Bigr),
\]
a unit vector of $\Q(\sqrt d)^2$, with mirror image $\bar u_n$. Rotating a vector $v$ by $\gamma\in G_N$ gives a
unit vector $\gamma v$ of $\Q(\sqrt d)^2$.

\begin{theorem}\label{thm:sets}
Let $d\equiv11\pmod{12}$, $N=5^k$, and let $U$ be the following set of unit vectors.
\begin{enumerate}
\item If $d\equiv23\pmod{24}$ and $d<21\cdot5^{k-1}$ (for example $5^k\ge d/4$): $U=G_N\cup G_Nu_1\cup G_N\bar u_1$.
\item If $d\equiv11\pmod{24}$, $s=v_3(d+1)$, $n$ is a positive integer with $n\equiv3\pmod4$ and
$n\equiv1\pmod{3^{s+1}}$ (for example
$n=1+2\cdot3^{s+1}$), and $5^k>\frac{2\sqrt{10}}3\,n(n+d)$: $U=G_N\cup G_Nu_1\cup G_N\bar u_1\cup G_Nu_n$.
\end{enumerate}
Then no homomorphism $\xi$ from the group $\Z U$ generated by $U$ to $\R/\Z$ maps every vector of $U$ into the arc
$[1/3,2/3]$.
\end{theorem}

\begin{proof}[Proof of Theorem~\ref{thm:main}]
By \cite[Theorem~1]{Winding}, a Cayley graph $\Cay(\Gamma,S)$ of an abelian group $\Gamma$, with $S=-S$, is
$3$-colourable if and only if some character $\xi\colon\Gamma\to\R/\Z$ maps $S$ into $[1/3,2/3]$. By
Theorem~\ref{thm:sets}, $\Cay(\Z U,\pm U)$ is not $3$-colourable. It is a subgraph of the plane over $\Q(\sqrt d)$.
(If $d$ is not squarefree, its squarefree part is again $\equiv11\pmod{12}$, and the field is the same.)
\end{proof}

By the theorem of de Bruijn and Erd\H{o}s \cite{dBE}, each of these planes contains a finite unit-distance graph with no
proper $3$-colouring, and deleting vertices one at a time gives one with chromatic number exactly $4$; it is
triangle-free, since $\sqrt3\notin\Q(\sqrt d)$. We do not exhibit it: the
direction ``$3$-colourable $\Rightarrow$ character'' of \cite[Theorem~1]{Winding} averages over the group with an
invariant mean and is not constructive. For $27$ of these fields, explicit vertex-critical graphs with $71$ to
$1\,404$ vertices and no proper $3$-colouring are given in \cite{QP}, with proofs that do not use characters.

Section~\ref{sec:chars} turns characters into complex numbers, Section~\ref{sec:shape} describes the set $S_N$ of
characters of $\frac1N\Z[i]$ that are allowed on $G_N$, and Sections~\ref{sec:23} and~\ref{sec:11} prove the two
cases of Theorem~\ref{thm:sets}. The bound on $k$ is generous: for $d=11$ the theorem needs $5^k=3125$, but $5^k=25$
(40 vectors up to sign) already works, as an exact certificate checked by computer shows (Section~\ref{sec:checks}).

\section{Characters as complex numbers}\label{sec:chars}

Identify $\Q^2$ with $\Q(i)$, and write a vector of $\Q(\sqrt d)^2$ as $X+\sqrt dY$ with $X,Y\in\Q(i)$: the vector
$(X_1+\sqrt dY_1,\,X_2+\sqrt dY_2)$ corresponds to $X=X_1+iX_2$, $Y=Y_1+iY_2$. For $\gamma\in\Q(i)$ with $|\gamma|=1$,
multiplying $X$ and $Y$ by $\gamma$ is the rotation by $\gamma$, an isometry of $\Q(\sqrt d)^2$. In this notation
$u_n=X+\sqrt dY$ with $X=(n^2-d)/(n^2+d)$ and $Y=2ni/(n^2+d)$.

Let $U$ be one of the sets of Theorem~\ref{thm:sets} and $T$ the least common multiple of the denominators $n^2+d$
of the vectors $u_n$ used. Every vector of $U$ lies in $\Gamma'=M+\sqrt dM$, $M=\frac1{NT}\Z[i]$, because
$N\gamma\in\Z[i]$ for $\gamma\in G_N$ and $u_n=X+\sqrt dY$ with $X,Y\in\frac1T\Z[i]$. Every character $\xi$ of $\Z U$
extends to $\Gamma'$, since $\R/\Z$ is divisible. The characters of $\Gamma'\cong\Z^4$ are
$X+\sqrt dY\mapsto\operatorname{Re}(\bar c_1X)+\operatorname{Re}(\bar c_2Y)\bmod 1$ with $c_1,c_2\in\C$ (on the basis
$\frac1{NT},\frac i{NT}$ of $M$ these take the arbitrary real values $\frac{\operatorname{Re}c}{NT}$,
$\frac{\operatorname{Im}c}{NT}$). For a
vector $v=X+\sqrt dY$ put $w_v=c_1\bar X+c_2\bar Y\in\C$. Then
\begin{equation}\label{eq:wv}
\xi(\gamma v)=\operatorname{Re}(\bar w_v\gamma)\bmod1\qquad(\gamma\in G_N),
\end{equation}
and $v\mapsto w_v$ is $\Q$-linear, so every linear relation with rational coefficients among vectors holds,
exactly in $\C$, among the corresponding numbers $w_v$. Let $I=[1/3,2/3]$ and
\[
S_N=\{c\in\C:\ \operatorname{Re}(\bar c\gamma)\in I+\Z\ \text{for every }\gamma\in G_N\}.
\]
If $\xi(U)\subseteq I+\Z$, then $w_v\in S_N$ for each $v\in\{1,u_1,\bar u_1,u_n\}$, where $1=(1,0)$. The set $S_N$
is invariant under $N\Z[i]$ (as $N\gamma\in\Z[i]$), under $c\mapsto-c$ and under multiplication by $i$.

\section{The shape of \texorpdfstring{$S_N$}{S\_N}}\label{sec:shape}

The definitions of $G_N$ and $S_N$ make sense for every $k\ge0$; $G_1=\{\pm1,\pm i\}$. Let $\rho=(3+4i)/5=(2+i)/(2-i)$,
$c^*_N=N(1+i)/2$, $Q_N=\{\frac N3(\alpha+\beta i):\alpha,\beta\in\{1,2\}\}$,
$\Sq_0=\{x+iy:|x|,|y|\le1/6\}$, and for $k\ge0$
\[
P_k=\{x\in\C:\ \bar x\rho^j\in\Sq_0\ \text{for all integers }|j|\le k\},
\]
a closed convex polygon; $P_0=\Sq_0$ and $P_k\subseteq P_1$ for $k\ge1$. Let $\Lambda_\pm=\frac{2\pm i}5\Z[i]$; both
contain $\Z[i]$ with index~$5$.

\begin{lemma}\label{lem:GN}
$G_N=\{\varepsilon\rho^j:\varepsilon\in\{\pm1,\pm i\},\ |j|\le k\}$, a set of $4(2k+1)$ elements.
\end{lemma}

\begin{proof}
Let $\gamma\in G_N$, $\gamma=\alpha/5^t$ with $\alpha\in\Z[i]$ and $t\le k$ least. Then $|\alpha|^2=25^t$, so
$\alpha=\varepsilon(2+i)^e(2-i)^f$ with $e+f=2t$; if $t\ge1$, minimality excludes $e,f\ge1$, so
$\gamma=\varepsilon(2\pm i)^{2t}/5^t=\varepsilon\rho^{\pm t}$. These elements are distinct, since $\rho^j$ has exact
denominator $5^{|j|}$.
\end{proof}

\begin{lemma}\label{lem:points}
\textup{(a)} $\bar c^*_N\gamma\in\frac{1+i}2+\Z[i]$ for every $\gamma\in G_N$. Hence $c^*_N+x\in S_N$ if and only if
$\bar x\gamma\in\Sq_0+\Z[i]$ for every $\gamma\in G_N$; in particular $c^*_N+P_k\subseteq S_N$.
\textup{(b)} $Q_N\subseteq S_N$.
\end{lemma}

\begin{proof}
(a) By Lemma~\ref{lem:GN}, $\bar c^*_N\gamma=\frac{5^{k-|j|}}2(1-i)\varepsilon\pi$ with $\pi=(2\pm i)^{2|j|}$. As
$\varepsilon\pi$ has odd norm, $\varepsilon\pi\equiv1\pmod{1+i}$ and $(1-i)\varepsilon\pi\equiv1-i\pmod2$, so both
coordinates are odd. Since $I=\frac12+[-\frac16,\frac16]$, the conditions for $\gamma$ and $i\gamma$ together say
$\bar x\gamma\in\Sq_0+\Z[i]$.
(b) For $q=\frac N3(\alpha+\beta i)$, $\bar q\gamma=\frac{5^{k-|j|}}3(\alpha-\beta i)\varepsilon\pi$. In
$\Z[i]/3=\F_9$, the elements with both coordinates nonzero, $\pm1\pm i$, are the non-squares, since $1+i$
generates $\F_9^\times$ and $(1+i)\mu_4=\{\pm1\pm i\}$. Now $\alpha-\beta i$ is a non-square and $\varepsilon\pi$ a
nonzero square, so both coordinates of $(\alpha-\beta i)\varepsilon\pi$ are prime to $3$, and those of $\bar q\gamma$
lie in $\{\frac13,\frac23\}+\Z$.
\end{proof}

\begin{lemma}\label{lem:numbers}
\textup{(a)} $\max_{x\in P_1}|x|=\sqrt{10}/18\approx0.1757$.
\textup{(b)} Every point of $(\Lambda_+\cup\Lambda_-)\setminus\Z[i]$ is at distance at least
$\sqrt2/6\approx0.2357$ from $\Sq_0+\Z[i]$, and none lies in $[-\frac13,\frac13]^2+\Z[i]$.
\end{lemma}

\begin{proof}
(a) $P_1$ is cut out by the six slabs $|\langle x,\nu\rangle|\le\frac16$, $\nu\in\{1,i,\rho,i\rho,\bar\rho,i\bar\rho\}$,
whose directions modulo $180^\circ$ are $0^\circ$, $\theta$, $90^\circ-\theta$, $90^\circ$, $90^\circ+\theta$,
$180^\circ-\theta$ with $\theta=\arccos\frac45$. So $P_1$ is a $12$-gon circumscribed about the circle of radius
$\frac16$; consecutive normals make angles $\theta$ (four times) and $90^\circ-2\theta$ (twice), and the vertices
are at distance $\frac16/\cos\frac\theta2=\frac16\cdot\frac{\sqrt{10}}3$, for example $(\frac16,\frac1{18})$, or
closer.
(b) The nonzero classes of $\Lambda_+/\Z[i]$ are those of $\pm\frac{2+i}5$ and $\pm\frac{-1+2i}5$, whose
representatives in $(-\frac12,\frac12)^2$ have one coordinate $\pm\frac25$ and the other $\pm\frac15$. The point of
$\Sq_0$ nearest to $\frac25+\frac i5$ is $\frac16+\frac i6$, at distance $\sqrt{(7/30)^2+(1/30)^2}=\sqrt2/6$; other
translates of $\Sq_0$ are farther. The coordinate $\pm\frac25$ is at distance $\frac25>\frac13$ from $\Z$.
$\Lambda_-$ is the mirror image of $\Lambda_+$.
\end{proof}

\begin{lemma}[unique lifts]\label{lem:lifts}
Let $k\ge1$, $N=5^k$, $N'=5^{k-1}$. If $c,c'\in S_N$ and $c'-c\in N'\Z[i]$, then $c'-c\in N\Z[i]$.
\end{lemma}

\begin{proof}
Write $c'=c+N'\lambda$. Then $\bar c'\rho^k-\bar c\rho^k=N'\bar\lambda\rho^k=\bar\lambda(2+i)^{2k-1}/(2-i)\in\Lambda_+$.
Both $\bar c\rho^k$ and $\bar c'\rho^k$ lie in $(I+\Z)+i(I+\Z)$ (the conditions for $\rho^k$ and $i\rho^k$), so the
difference lies in $[-\frac13,\frac13]^2+\Z[i]$, hence in $\Z[i]$ by Lemma~\ref{lem:numbers}(b): that is,
$(2-i)\mid\bar\lambda$, or $(2+i)\mid\lambda$. With $\rho^{-k}$ and $\Lambda_-$, $(2-i)\mid\lambda$. So $5\mid\lambda$.
\end{proof}

\begin{proposition}\label{prop:shape}
For every $k\ge1$, $S_{5^k}=(c^*_N+P_k+N\Z[i])\cup(Q_N+N\Z[i])$. So every $c\in S_N$ is $c=e+x$ with either
$e\in c^*_N+N\Z[i]$ and $|x|\le\sqrt{10}/18$, or $e\in Q_N+N\Z[i]$ and $x=0$.
\end{proposition}

\begin{proof}
``$\supseteq$'' is Lemma~\ref{lem:points}. For ``$\subseteq$'' we induct on $k$; write $N'=5^{k-1}$ and note that
$S_N\subseteq S_{N'}$ and $c^*_N-c^*_{N'}=2N'(1+i)\in N'\Z[i]$.

\emph{$k=1$.} $S_1=\frac{1+i}2+\Sq_0+\Z[i]$, so $c\in S_5$ is $c=c^*_5+x+\lambda$ with $x\in\Sq_0$, $\lambda\in\Z[i]$.
By Lemma~\ref{lem:points}(a), the conditions at $\rho$ say $\bar x\rho+\bar\lambda\rho\in\Sq_0+\Z[i]$, where
$\bar\lambda\rho\in\Lambda_+$. If $\bar\lambda\rho\notin\Z[i]$, Lemma~\ref{lem:numbers}(b) gives
$|x|=|\bar x\rho|\ge\sqrt2/6$, and as $x\in\Sq_0$, $x$ is a corner $(\pm1\pm i)/6$; likewise with $\bar\rho$. If $x$
is not a corner, then $(2+i)\mid\lambda$ and $(2-i)\mid\lambda$, so $c\equiv c^*_5+x\pmod{5\Z[i]}$, and
$\bar x\rho^{\pm1}\in\Sq_0+\Z[i]$ with $|x|<\frac56$ gives $\bar x\rho^{\pm1}\in\Sq_0$: $x\in P_1$. If $x$ is a
corner, $c\equiv\frac{1+i}2+\frac{\pm1\pm i}6=\frac{\alpha+\beta i}3\pmod{\Z[i]}$ with $\alpha,\beta\in\{1,2\}$; the
point $q=\frac53(\alpha'+\beta'i)\in Q_5$ with $\alpha'\equiv2\alpha$, $\beta'\equiv2\beta\pmod3$ lies in $S_5$ and
$q\equiv c\pmod{\Z[i]}$, so $c\in q+5\Z[i]$ by Lemma~\ref{lem:lifts}.

\emph{$k\ge2$.} Let $c\in S_N\subseteq S_{N'}$. If $c\equiv\frac{N'}3(\alpha+\beta i)\pmod{N'\Z[i]}$, then
$q=\frac N3(\alpha'+\beta'i)\in Q_N$ with $\alpha'\equiv2\alpha$, $\beta'\equiv2\beta\pmod3$ satisfies
$q\equiv c\pmod{N'\Z[i]}$ and $q\in S_N$, so $c\in q+N\Z[i]$ by Lemma~\ref{lem:lifts}. Otherwise
$c=c^*_N+x+N'\lambda$ with $x\in P_{k-1}\subseteq P_1$ and $\lambda\in\Z[i]$, and the conditions at $\rho^k$ say
$\bar x\rho^k+\mu\in\Sq_0+\Z[i]$ with $\mu=N'\bar\lambda\rho^k\in\Lambda_+$. If $\mu\notin\Z[i]$,
Lemma~\ref{lem:numbers} gives $\sqrt2/6\le|x|\le\sqrt{10}/18$, which is impossible. So $(2+i)\mid\lambda$, likewise
$(2-i)\mid\lambda$, hence $c\equiv c^*_N+x\pmod{N\Z[i]}$, and $\bar x\rho^{\pm k}\in\Sq_0$, so $x\in P_k$.
\end{proof}

\emph{Types.} By Proposition~\ref{prop:shape} write each $c\in S_N$ as $c=N\varepsilon+x$ with
$|x|\le r:=\sqrt{10}/18$ and $\varepsilon$ of \emph{type~c}, $\varepsilon\in\frac{1+i}2+\Z[i]$, or of \emph{type~q},
$\varepsilon\in\frac{\alpha+\beta i}3+\Z[i]$ with $\alpha,\beta\in\{1,2\}$; in both cases
$N\varepsilon\in\frac N6\Z[i]$. Let $\Z_{(p)}[i]$ be the elements of $\Q(i)$ whose denominators are prime to $p$,
and $v_3$ the $3$-adic valuation of $\Q(i)$ ($3$ is inert in $\Z[i]$). Then: $\varepsilon$ has type~q if and only
if $v_3(\varepsilon)=-1$; type~c if and only if $\varepsilon\notin\Z_{(2)}[i]$, and then
$\varepsilon\equiv h:=\frac{1+i}2\pmod{\Z_{(2)}[i]}$; $mh\equiv h$ for odd $m$, and $2h\in\Z[i]$.

\begin{lemma}[exact relations]\label{lem:exact}
Let $\sigma_1,\dots,\sigma_\ell\in S_N$, $\sigma_j=N\varepsilon_j+x_j$, and let $m_1,\dots,m_\ell\in\Z$ with
$\sum m_j\sigma_j=0$. If
$N>6r\sum|m_j|$, then $\sum m_j\varepsilon_j=0$.
\end{lemma}

\begin{proof}
$\sum m_jN\varepsilon_j=-\sum m_jx_j$; the left side lies in $\frac N6\Z[i]$, whose nonzero elements have modulus
at least $\frac N6$, and the right side has modulus at most $r\sum|m_j|<\frac N6$.
\end{proof}

\section{The case \texorpdfstring{$d\equiv23\pmod{24}$}{d = 23 mod 24}}\label{sec:23}

Let $a=\frac{d+1}4$ and $b=\frac{d-1}2$: here $a\equiv0\pmod6$, $b$ is odd and $3\nmid b$. The relation
$a(u_1+\bar u_1)+b\cdot1=0$ gives, for $z=w_1$, $w=w_{u_1}$, $w'=w_{\bar u_1}$ in $S_N$,
\[
a(w+w')+bz=0.
\]
By Proposition~\ref{prop:shape} write $w=e_1+x_1$, $w'=e_2+x_2$, $z=e_3+y$, with
$e_j\in(c^*_N+N\Z[i])\cup(Q_N+N\Z[i])$, $x_1,x_2,y\in P_k$, and the corresponding $x_j$ or $y$ equal to $0$ when
$e_j\in Q_N+N\Z[i]$. Then
\[
a(e_1+e_2)+be_3=-(ax_1+ax_2+by).
\]
As $a$ is even and $3\mid a$, $ac^*_N=\frac a2N(1+i)$ and $a\frac N3(\alpha+\beta i)$ lie in $N\Z[i]$; as $b$ is odd
and $3\nmid b$, $be_3$ lies in $c^*_N+N\Z[i]$ or in $Q_N+N\Z[i]$ with $e_3$. So the left side $u+iv$ lies in
$c^*_N+N\Z[i]$, where $|u|,|v|\ge\frac N2$, or in $Q_N+N\Z[i]$ ($-Q_N\equiv Q_N$), where $|u|,|v|\ge\frac N3$. As
$P_k$ is convex and symmetric, the right side lies in $(2a+b)P_k=dP_k$ in the first case and in
$2aP_k=\frac{d+1}2P_k$ in the second ($y=0$). On $P_1\supseteq P_k$ the conditions at $\rho$ and $\bar\rho$ give
$|3u\pm4v|\le\frac56$, so $3|u|+4|v|\le\frac56$. Hence the first case needs $\frac{7N}2\le\frac{5d}6$, that is
$d\ge\frac{21N}5$, and the second needs $\frac{7N}3\le\frac{5(d+1)}{12}$, that is $d+1\ge\frac{28N}5$. If
$d<21N/5$, neither holds. This proves Theorem~\ref{thm:sets}(1); and $5^k\ge d/4$ gives $d\le4N<21N/5$.

\begin{remark}
For $k\le5$ the bound is sharp: if $d\ge21N/5$, then $c^*_N/d\in P_k$, since for $k\le5$ the conditions at
$\rho^{\pm1}$ are the binding ones, and this gives a character with all three types~c. An exact computation (cells of
$S_N$, Minkowski sums and the lattice points of $N\Z[i]$, with a second method as a cross-check and an explicit
character in every feasible case) confirms, for every $d\equiv23\pmod{24}$ up to $1295$ when $N=5,25,125$ and up to
$2999$ when $N=625$, that the set of Theorem~\ref{thm:sets}(1) admits no such character exactly when $d<21N/5$. For
$k\ge6$ the bound is not sharp: as $(1-i)\rho^6=(1457-22049i)/15625$, the point $c^*_N/d$ lies in $P_6$ only when
$d\ge66147N/15625$, and for $N=5^6$ an exact decision (all choices of types and all lattice points) finds no
character for the $22$ values $65\,639\le d\le66\,143$ with $d\equiv23\pmod{24}$, although $d\ge21N/5=65\,625$; the
first value with a character is $66\,167$.
\end{remark}

\section{The case \texorpdfstring{$d\equiv11\pmod{24}$}{d = 11 mod 24}}\label{sec:11}

Now $a$ is odd, and the argument fails: $a+b$ is even and $3\mid a$, so $a(c^*_N+q)+bc^*_N\in N\Z[i]$ for $q\in Q_N$,
and since $\gcd(a,b)=1$ suitable lifts $z=c^*_N+N\lambda_3$, $w=c^*_N+N\lambda_1$, $w'=q+N\lambda_2$ solve
$a(w+w')+bz=0$ exactly. So $G_N\cup G_Nu_1\cup G_N\bar u_1$ admits such a character for every $k$. The vector $u_n$ excludes it. Comparing
$X$ and $Y$ gives $(n^2+d)u_n=(n-1)(n+d)\cdot1+n(1+d)u_1$, so with $y=w_{u_n}$,
\[
a(w+w')+bz=0,\qquad(n^2+d)y-(n-1)(n+d)z-n(1+d)w=0.
\]
As $n\ge1$, $(n-1)(n+d)$ and $n(1+d)$ are nonnegative, and the sums of their absolute values are $2a+b=d$ and
$2n(n+d)\ge d$. (For negative $n$ the second sum is larger, and the hypothesis on $N$ would have to change.) As
$N>\frac{2\sqrt{10}}3n(n+d)=12r\,n(n+d)$, Lemma~\ref{lem:exact} applies to both relations: writing
$z=N\varepsilon_z+x_z$ and so on,
\begin{equation}\label{eq:exact11}
a(\varepsilon_w+\varepsilon_{w'})=-b\varepsilon_z,\qquad
(n^2+d)\varepsilon_y=(n-1)(n+d)\varepsilon_z+n(1+d)\varepsilon_w,
\end{equation}
exactly in $\Q(i)$.

\emph{Step 1: $\varepsilon_z$ has type~c and $v_3(\varepsilon_z)\ge s-1$.} We have $v_3(a)=s\ge1$, $3\nmid b$, and
$v_3(\varepsilon_w+\varepsilon_{w'})\ge-1$, so $v_3(\varepsilon_z)\ge s-1\ge0$.

\emph{Step 2 (at $2$): $\varepsilon_y$ has type~c if and only if $\varepsilon_w$ has type~q.} As $d\equiv3\pmod8$
and $n\equiv3\pmod4$, each of $n^2+d$, $(n-1)(n+d)$ and $n(1+d)$ is $4$ times an odd integer; write the second
relation as $A'\varepsilon_y=B'\varepsilon_z+C'\varepsilon_w$ with $A',B',C'$ odd. Modulo $\Z_{(2)}[i]$,
$B'\varepsilon_z\equiv h$, and $C'\varepsilon_w\equiv h$ or $0$ according as $\varepsilon_w$ has type~c or~q. So
$A'\varepsilon_y\in\Z_{(2)}[i]$ exactly when $\varepsilon_w$ has type~c.

\emph{Step 3 (at $3$): $\varepsilon_y$ has type~q if and only if $\varepsilon_w$ has type~q.} Let
$e=v_3(n-1)\ge s+1$; $3\nmid n(n+1)$. Then $n^2+d=(n-1)(n+1)+(d+1)$ and $n+d=(n-1)+(d+1)$ have valuation $s$, so do
$n(1+d)$, and $v_3((n-1)(n+d))=e+s$. Dividing the second relation by $3^s$ gives
$A''\varepsilon_y=B''\varepsilon_z+C''\varepsilon_w$ with $3\nmid A''C''$ and
$v_3(B''\varepsilon_z)\ge e+s-1\ge2s\ge2$. So $v_3(\varepsilon_y)=-1$ if and only if $v_3(\varepsilon_w)=-1$.

Steps 2 and 3 contradict each other whatever the type of $\varepsilon_w$. This proves Theorem~\ref{thm:sets}(2).

\begin{remark}
The argument of Section~\ref{sec:23} does not work with one vector in place of $u_1$ and $u_n$: if
$v=(\mu+i\nu\sqrt d)/\tau$ with $\mu^2+d\nu^2=\tau^2$ in lowest terms and $d\equiv3\pmod8$, then $4\nmid\tau$ (if $\tau$ is even, $\mu,\nu$ are odd and $\mu^2+d\nu^2\equiv4\pmod8$), while the
argument of Section~\ref{sec:23} needs the coefficient $\tau/\gcd(\tau,2\mu)$ of the relation $\tau(v+\bar v)=2\mu\cdot1$
to be even. Correspondingly, $x^2+y^2$ is anisotropic over $\Q_2(\sqrt d)$ for $d\equiv3\pmod8$ and isotropic for
$d\equiv7\pmod8$.
\end{remark}

\section{Checks}\label{sec:checks}

Every computation that the proofs use is checked in exact arithmetic by a program in \cite{Repo}: Lemmas
\ref{lem:GN} and \ref{lem:points} for $k\le6$, both numbers of Lemma~\ref{lem:numbers}, and $S_{5^k}$ itself for
$k=1,2,3$ by exact polygon clipping (a $12$-, $20$- and $28$-gon around $c^*_N$ of areas $17/189$, $61/693$ and
$437/4977$, and the four points of $Q_N$), independently of the induction; a referee's own program confirmed
Proposition~\ref{prop:shape} with equality for $k\le4$, and found no exception among $10^7$ sampled points for
$k=1,2$. Exact certificates (branch and bound over the integer values of a basis of the relations, with a rational
Farkas vector at every leaf), each accepted by two checkers written separately (\nolinkurl{check_w.py} and
\nolinkurl{check_w_indep.py}), show that the sets of Theorem~\ref{thm:sets} have no such character: in case~(1),
with the $N$ that the proof needs, $d=23,47,71,95$ with $N=25$ and $d=119,143,\dots,311$ ($d\equiv23\pmod{24}$,
squarefree) with $N=125$, an independent check of Section~\ref{sec:23}; in case~(2), already with a smaller $N$
than the proof needs, $d=11,35,59$ with $N=25$ and $d=83,131,155$ with $N=125$. A floating-point solver confirms case~(2) with the $N$ that the proof needs for
$d=11,35,59,83,131,179$. Finally, for a finite set $\Delta$ of unit vectors, Proposition~\ref{prop:shape} and
Lemma~\ref{lem:exact} reduce the question whether $G_N(\Delta\cup\{1\})$ admits such a character for all large $N$ to
the existence of a character with values in $\{\frac13,\frac12,\frac23\}$, a computation in $(\Z/6)^4$; it confirms
Theorem~\ref{thm:sets}(2) for all $191$ squarefree $d\equiv11\pmod{24}$ below $5000$, and shows that for $d=11$ no
two vectors $(\mu+i\nu\sqrt{11})/\tau$ with $\nu\le7$, $\tau\le60$ suffice. Independently of the proof, a computer
search over denominators found a certificate of the same kind, for other sets of unit vectors, for each of the
$155$ squarefree $d\equiv11\pmod{12}$ below $2000$ and for $593$ of the $611$ between $2000$ and $10\,000$; both
checkers accept all of them, except two of more than $50$ MB, which were checked once and not kept
(\nolinkurl{data/quadratic_planes/winding/scan} in \cite{Repo}).

\emph{Formal proof.} Theorem~\ref{thm:main} and Corollary~\ref{cor:padic} are also proved in Lean~4 with Mathlib
(\nolinkurl{lean/FourColours.lean} and \nolinkurl{lean/PadicFour.lean} in \cite{Repo}, the values
$\chi(\Q_2^2)=2$ and $\chi(\Q_3^2)=3$ in \nolinkurl{lean/PadicPlanes.lean}, and Theorem~W for an arbitrary abelian
group and a finite set $S$ in \nolinkurl{lean/TheoremWInf.lean}, by averages over boxes and a limit along an
ultrafilter); the files depend only on Lean's three standard axioms. The formal proof uses a cruder form of
Proposition~\ref{prop:shape}, which suffices for Theorem~\ref{thm:main} but not for the bound
$d<21\cdot5^{k-1}$: every $c\in S_N$ is $N(\frac{1+i}2+t)+x$ with $t\in\Z[i]$ and $|x|^2<\frac1{18}$, or
$N(\frac{\alpha+\beta i}3+t)$ with $\alpha,\beta\in\{1,2\}$. This follows by the induction of
Proposition~\ref{prop:shape} started at $k=0$, where $S_1=([\frac13,\frac23]+\Z)^2$ and the corners of the square are
the points of type $q$; in the step, $\frac1{18}=(\frac7{30})^2+(\frac1{30})^2$ replaces the numbers of
Lemma~\ref{lem:numbers}, and the points of type $q$ are handled by a congruence modulo $15$. The formal analogue of Lemma~\ref{lem:exact}
assumes $N>2\sum_j|m_j|$, for values written as $\frac N6E+x$ with $E$ an integer and $|x|\le\frac14$.

\section*{Acknowledgements}
The author thanks the authors of the works cited, whose questions and results this note builds on.

\begin{thebibliography}{99}
\bibitem{Davies} J. Davies, \emph{Chromatic number of spacetime}, arXiv:2308.16885 (2023, revised 2024).
\bibitem{dBE} N. G. de Bruijn and P. Erd\H{o}s, \emph{A colour problem for infinite graphs and a problem in the
theory of relations}, Nederl. Akad. Wetensch. Proc. Ser. A \textbf{54} = Indag. Math. \textbf{13} (1951),
371--373.
\bibitem{deGrey} A. D. N. J. de Grey, \emph{The chromatic number of the plane is at least 5}, Geombinatorics
\textbf{28} (2018), 18--31; arXiv:1804.02385.
\bibitem{EI} G. Exoo and D. Ismailescu, \emph{The chromatic number of the plane is at least $5$: a new proof},
Discrete Comput. Geom. \textbf{64} (2020), 216--226; arXiv:1805.00157.
\bibitem{Fischer1990} K. G. Fischer, \emph{Additive $K$-colorable extensions of the rational plane}, Discrete
Math. \textbf{82} (1990), 181--195.
\bibitem{Johnson1987} P. D. Johnson Jr., \emph{Two-colorings of real quadratic extensions of $\Q^2$ that forbid
many distances}, Congr. Numer. \textbf{60} (1987), 51--58.
\bibitem{Madore} D. A. Madore, \emph{The Hadwiger--Nelson problem over certain fields}, arXiv:1509.07023 (2015).
\bibitem{Moorhouse} G. E. Moorhouse, \emph{On the chromatic numbers of planes}, preliminary draft, revised
3 March 2010, \url{https://www.ericmoorhouse.org/pub/chromatic.pdf}.
\bibitem{MoorhouseTalk} G. E. Moorhouse, \emph{Colouring the plane}, talk at Designs, Codes \& Geometries
2010, slides at \url{https://www.ericmoorhouse.org/slides/ebertfest.pdf}.
\bibitem{OpenAI} OpenAI, \emph{The Euclidean plane is not five-colorable}, preprint (23 September 2026), with a
formalization in Lean, \url{https://github.com/openai/math}.
\bibitem{PA} S. Gal\'an, \emph{Colouring the $p$-adic plane}, draft, 2026, \texttt{papers/padic-planes/} in
\cite{Repo}.
\bibitem{QP} S. Gal\'an, \emph{Real quadratic planes that need four colours}, draft, 2026,
\texttt{papers/quadratic-planes/} in \cite{Repo}.
\bibitem{Repo} S. Gal\'an, \emph{The Hadwiger--Nelson problem over number fields}, code, data and formal
proofs, 2026. Archived at Zenodo, \href{https://doi.org/10.5281/zenodo.22976635}{doi:10.5281/zenodo.22976635}
(all versions); source at \url{https://github.com/decalion89/chromatic-number-of-the-plane}; a fuller account is
\texttt{notes/four\_colours\_11\_mod\_12.md} there, and the programs are in
\texttt{data/quadratic\_planes/winding/family/}. These files and the Lean files cited here are on the branch
\texttt{claude/quadratic-planes} (pull request~35), not yet in a release.
\bibitem{Soifer} A. Soifer, \emph{The mathematical coloring book}, Springer, New York, 2009.
\bibitem{Winding} S. Gal\'an, \emph{Circular colourings of abelian Cayley graphs below four come from
characters}, draft, 2026, \texttt{papers/winding/} in \cite{Repo}.
\bibitem{Woodall} D. R. Woodall, \emph{Distances realized by sets covering the plane}, J. Combin. Theory Ser. A
\textbf{14} (1973), 187--200.
\end{thebibliography}

\end{document}
